% Generated by Listas/make_statements.py from pratica-lecture-03.md, item 13. Do not edit.
\textbf{Source.} Tadelis (2013), \emph{Game Theory: An Introduction}, Exercise 8.11 ``Resident Assistant Selection with a Twist'', pp.~172--173 {[}PDF 191--192{]}. Reproduced verbatim. Answer in the source: none in the book (the preface, p.~xiii, says about half are solved on Princeton UP's website; not checked).



\begin{quote}

\textbf{Resident Assistant Selection with a Twist:} Two staff managers in the ΠΒΦ sorority, the house manager (player 1) and kitchen manager (player 2), must select a resident assistant from a pool of three candidates: \(\{a, b, c\}\). Player 1 prefers \(a\) to \(b\), and \(b\) to \(c\). Player 2 prefers \(b\) to \(a\), and \(a\) to \(c\). The process that is imposed on them is as follows: First, the house manager vetoes one of the candidates and announces the veto to the central office for staff selection and to the kitchen manager. Next the kitchen manager vetoes one of the remaining two candidates and announces it to the central office. Finally the director of the central office assigns the remaining candidate to be a resident assistant at ΠΒΦ.



\begin{enumerate}

\def\labelenumi{\alph{enumi}.}

\item

  Model this as an extensive-form game (using a game tree) in which a player's most-preferred candidate gives a payoff of 2, the second most-preferred candidate gives a payoff of 1, and the last candidate gives 0.

\item

  Find the subgame-perfect equilibrium for this game. Is it unique?

\item

  Are there Nash equilibria that are not subgame-perfect equilibria?

\item

  Now assume that before the two players play the game, player 2 can send an alienating e-mail to one of the candidates, which would result in that candidate withdrawing her application. Would player 2 choose to do this, and if so, with which candidate?

\end{enumerate}

\end{quote}



\begin{center}\rule{0.5\linewidth}{0.5pt}\end{center}

